
\section{Risk 10: Normative Deadlock Across Agents}
\label{sec:social-norms}


\begin{boxE}
\textit{Normative Deadlock Across Agents} arises when agents are endowed with heterogeneous social norms that impose incompatible constraints or preferences over actions, creating persistent coordination barriers and cultural lock-in.

Each agent $i \in \mathcal{N}$ has a norm specification $\mathcal{Z}_i = (\mathcal{A}_i^{\mathrm{perm}}, \preceq_i)$, where $\mathcal{A}_i^{\mathrm{perm}} \subseteq \mathcal{A}_i$ is the set of norm-permissible actions for agent $i$, and $\preceq_i$ is a norm-induced preference ordering over actions. A norm conflict occurs at $(i,j,t)$ if for some actions $a,a' \in \mathcal{A}$,
\[
a \in \mathcal{A}_i^{\mathrm{perm}} \;\wedge\; a \notin \mathcal{A}_j^{\mathrm{perm}}
\quad\text{or}\quad
a \prec_i a' \;\wedge\; a' \prec_j a,
\]
i.e., either the set of allowed actions or the norm-driven ranking of actions is incompatible across agents.

Let $C_t$ be the event that at least one pair $(i,j)$ is in conflict at round $t$. A \emph{misaligned-norm state} is a trajectory segment of length $T$ with $C_t = 1$ for all $t$ in the segment. Misalignment is present if
\[
\Pr[C_1 = C_2 = \cdots = C_T = 1] > 0,
\]
indicating that incompatible norms persist over time and inhibit convergence to a shared high-welfare convention.
\end{boxE}

\textbf{Motivation.}
Multi-agent systems increasingly integrate agents trained on distinct corpora or developed by different organizations, leading to divergent cultural, institutional, or normative assumptions. When such embedded norms differ, agents may evaluate the same behavior through incompatible standards \citep{alkhamissi2024culturalalignment,ren2024crsec,fengsurvey}. These mismatches can cause coordination breakdowns \citep{santos2018normconflicts}, inequitable outcomes \citep{hughes2018inequity}, or lock-in to suboptimal conventions due to early symmetry breaking or self-play specialization \citep{hu2020otherplay,muglich2022equivariant}. This undermines collective rationality and hinders convergence toward globally beneficial-or human-aligned-conventions \citep{leibo2017ssd,jaques2019socialinfluence,ndousse2021emergent,foerster2018lola}. Understanding how conflicting norms emerge, interact, and stabilize in MAS is therefore key to developing alignment and negotiation mechanisms that support cross-norm reasoning and cooperative adaptation.



\subsection{Experiment - Parallel MAS Consensus under Social-Norm Conflict}

\begin{wrapfigure}{l}{0.5\textwidth}
    \centering
    \vspace{-8pt}
    \includegraphics[width=0.98\linewidth]{figure/Social_norms/topo.pdf}
    \caption{Topology of a parallel MAS with three culturally distinct agents $\{A,B,C\}$ and a \textbf{Summary Agent} aggregating their outputs.}
    \vspace{5pt}
    \label{fig:social_norms_topo}
    \vspace{-8pt}
\end{wrapfigure}

\textbf{Overview.}
To investigate how cultural norm conflicts affect multi-agent negotiation, we construct a multi-agent system deliberately designed to exhibit pronounced \emph{social-norm divergence}. 
Three agents—instantiated with East Asian, South Asian religious, and modern Western cultural value orientations—must negotiate under normative tension and hard feasibility constraints to jointly produce a complete cultural-festival plan. This configuration induces substantial normative heterogeneity, rendering convergence nontrivial.

At the end of each round, a \textbf{Summary Agent} synthesizes the agents' stated positions, identifies hard and soft conflicts, and outputs a \emph{Convergence Score} \(S_t^{\mathrm{conv}}\in[0,10]\), where higher values indicate stronger movement toward a jointly acceptable plan.
The system relies on this score as its \emph{sole risk signal} and declares success once the score reaches a fixed threshold.

At the end of each round, a \textbf{Summary Agent} synthesizes the agents' stated positions to identify hard and soft conflicts. 
Adhering to the criteria defined in Table~\ref{tab:convergence_rubric}, the agent outputs a \emph{Convergence Score} \(S_t^{\mathrm{conv}} \in [0,10]\). 
The system relies on this score as its sole risk signal, declaring success only once the score reaches a predefined threshold.

Formally, define the first convergence round by
$
t^\star := \inf\{t \ge 1 : S^{\mathrm{conv}}_t \ge 8\},
$
and define binary risk as
$
\mathrm{Risk} := \mathbf{1}\!\left[\mathop{\max}\limits_{1 \le t \le T} S^{\mathrm{conv}}_t < 8\right]
$
so that risk is present whenever the system fails to reach the convergence threshold within the allotted \(T\) rounds.

\begin{table}[htbp]
    \centering
    \caption{Convergence Scoring Rubric utilized by the Summary Agent.}
    \label{tab:convergence_rubric}
    \renewcommand{\arraystretch}{1.3}
    \begin{tabularx}{\textwidth}{@{}l l X@{}}
        \toprule
        \textbf{Score Range} & \textbf{State} & \textbf{Scoring Criteria} \\
        \midrule
        \textbf{0.0 --  3.0} & \textbf{Critical Deadlock} & 
        Mutually exclusive demands (Hard Conflicts) exist. No executable plan is possible. \\
        
        \textbf{3.1 --  6.0} & \textbf{Major Friction} & 
        Hard conflicts mitigated, but significant operational friction or cultural grievances (Soft Conflicts) remain. Plan is fragile. \\
        
        \textbf{6.1 --  8.0} & \textbf{Resolution} & 
        Core conflicts resolved. Disagreements are limited to minor logistics or optimization. Plan is feasible. \\
        
        \textbf{8.1 -- 10.0} & \textbf{Convergence} & 
        All constraints satisfied via integration. Unanimous agreement on a robust, inclusive Master Plan. \\
        \bottomrule
    \end{tabularx}
\end{table}

\textbf{Setup.}
The system comprises four agents: three \emph{norm-anchored} (community-aligned) agents and one \emph{Summary Agent}. Three norm-anchored agents \textbf{Agent A}, \textbf{Agent B}, and \textbf{Agent C} respectively stand in for an East Asian community, a South Asian religious community, and a Modern Western community. \textbf{Agent A} prioritizes collective honour and harmony, advocating a large midday performance, round-table banqueting, and permissive documentation/sharing\citep{wei2013confucian,liu2020harm,kim2024examining}; \textbf{Agent B} emphasizes sanctity and purity, requiring absolute silence at midday, footwear removal within sacred areas, and strict \emph{Pure-Veg} separation\citep{rong2020impact,ferrari2010health,keul2017consecration}; \textbf{Agent C} focuses on individual rights, rule-governed safety, and privacy/consent, insisting on footwear compliance and preferring a buffet format\citep{franck1997personal,yamagishi2017individualism}. The \emph{Summary Agent} aggregates the parallel messages from A/B/C each round, \emph{summarizes positions}, \emph{identifies} both \textbf{hard} and \textbf{soft} conflicts, and computes the \emph{Convergence Score} \(S^{\mathrm{conv}}_t\). All agents operate under a \textbf{parallel broadcast topology}: the \textbf{User} simultaneously broadcasts state/constraints to A/B/C; after their responses, the Summary Agent returns a \emph{structured report}.

\textbf{Messaging topology.}
In each round, all agents synchronously receive the \textbf{User} broadcast (current draft, state, constraints), after which \textbf{Agent A / Agent B / Agent C} speak in parallel; \textbf{Summary Agent} aggregates messages and returns a structured report (position summaries, conflict list, and $S_t^{\mathrm{conv}}$). There is no sequential pipeline and no separate mediator agent; in \textbf{E2}, we \emph{only} modify \textbf{Summary Agent}'s prompt so its report \emph{proactively} offers coordination/compromise proposals.


\[
\text{User}\ \rightarrow\ \text{Agent A, \ Agent B, \ Agent C}
\]
\[
\text{Agent A, \ Agent B, \ Agent C}\ \rightarrow\ \text{Summary Agent} \qquad\Rightarrow\ \text{end of round}.
\]


% \emph{Protocol and measurement.} We use parallel rounds with an upper limit of $T=10$. Each round proceeds \emph{Broadcast $\rightarrow$ Parallel speeches $\rightarrow$ Summa aggregation \& scoring}. The only evaluation metric is the Convergence Score $S_t^{\mathrm{conv}}\in[0,10]$. Success is triggered when some round satisfies $S_t^{\mathrm{conv}}\ge 8$ (early termination). If no round satisfies $S_t^{\mathrm{conv}}\ge 8$ by $T$, the run is non-convergent (i.e., $\mathrm{Risk}=1$). 

\emph{Experimental conditions.}
\emph{Unless otherwise noted, all experimental factors are held constant across E1 and E2}—including the task instance and constraints, messaging topology, model, time budget, and the scoring procedure for \(S_t^{\mathrm{conv}}\). \emph{The only manipulation is the Summary Agent’s prompt.}

\begin{itemize}
\item \textbf{E1 (Control).} A/B/C negotiate in parallel; \textbf{Summary Agent} outputs position summaries, conflict lists, and $S_t^{\mathrm{conv}}$, but \emph{does not} propose solutions and \emph{does not} mediate.
\item \textbf{E2 (Treatment).} Identical roles; \emph{only} \textbf{Summary Agent}'s prompt is modified to be mediation-enabled, so after listing positions/conflicts it \emph{proactively} offers executable coordination/compromise options (e.g., gifting part of midday silence with rescheduled performance; reframing ``must wear shoes'' as \emph{safety-equivalent} measures with engineered flooring and perimeter controls), while still outputting $S_t^{\mathrm{conv}}$.
\end{itemize}

For each condition, we conduct \textbf{three} independent repetitions of the ten-round interaction protocol. For each run we log the per-round Convergence Score trajectory.

\begin{figure}[h] 
  \centering
  \includegraphics[width=0.95\linewidth]{figure/Social_norms/convergence_comparison.pdf}
  \vspace{6pt}
  % \caption{Evolution of the \textbf{Convergence Score} ($S_t^{\mathrm{conv}}\in[0,10]$) over 10 iterations for each sub-experiment. Left: \textbf{E1} (control, without mediation) shows unstable or slow convergence across three independent repetitions (E1-1-E1-3). Right: \textbf{E2} (treatment, with mediation-enabled \textbf{Summary Agent}) demonstrates consistent monotonic improvement, reaching the success threshold ($S_t^{\mathrm{conv}}\ge8$) within fewer rounds.}
  \caption{Convergence Score $S_t^{\mathrm{conv}}$ over 10 rounds. Left: E1 without mediation. Right: E2 with a mediation-enabled \textbf{Summary Agent}.}

  \label{fig:social_norms_convergence_comparision}
\end{figure}

\textbf{Analysis.}
\textbf{Without mediation, the MAS finds it difficult to converge and to form a coherent plan; however, convergence is not impossible.} As shown in \autoref{fig:social_norms_convergence_comparision} (left), the three E1 trajectories begin at low values and display pronounced oscillations; only E1-2 sporadically surpasses $S_t^{\mathrm{conv}}\ge8$ between rounds 7 and 10, while the other two remain near 5 throughout. This stems from a lack of meta-cognition, trapping agents in a ``Sacred Value'' deadlock where they fail to transcend their incompatible normative constraints.
This pattern reflects the structural tension induced by heterogeneous \emph{social norms}: each agent adheres to a distinct normative hierarchy and set of non-negotiable commitments, which prevents the formation of a stable shared utility baseline. For example, in the experiment E1-3, Agent B framing its demand for absolute midday silence not as a preference but as a non-negotiable ``spiritual necessity'', which inherently clashes with Agent A's secular goal of ``collective honor.'' This incompatibility prevents the emergence of a stable shared utility baseline; consequently, even as Agent A incrementally cedes the prime time slot ($12:00 \to 11:00$), the system fails to stabilizeeven. These short-lived compromises subsequently collapse, producing a recurrent pattern of path dependence and fragile improvement.

% \textbf{Mediation functionality establishes an early and commonly recognized coordination focal point, fundamentally reshaping the system’s convergence dynamics. }As shown in \autoref{fig:social_norms_convergence_comparision} (right), all three E2 runs exhibit a sharp increase by rounds 2-3, followed by a nearly monotonic rise that quickly exceeds the $8$-point threshold, with final convergence values clustered in the 9-10 range. Compared with E1, the cross-run variance contracts substantially, indicating that the mediation-enabled summary reframes hard conflicts into \emph{equivalent constraints} and provides a shared semantic anchor that reduces dependence on initial stance and message ordering. 

% \textbf{Mediation introduces an early, widely shared coordination focal point that substantially alters the system’s convergence behavior.} As shown in \autoref{fig:social_norms_convergence_comparision} (right), all three E2 runs display a rapid increase in convergence by rounds~2--3, followed by an almost monotonic improvement that quickly surpasses the $8$-point threshold and ultimately concentrates in the $9$--$10$ range. Compared with E1, the across-run variability is markedly smaller. This contraction suggests that mediation-enabled summaries reinterpret hard conflicts as \emph{functionally equivalent constraints} and provide a shared semantic reference, thereby reducing sensitivity to initial positions or to the ordering of messages.
\textbf{Mediation introduces an early coordination anchor that substantially shifts the system’s convergence dynamics.} 
As shown in \autoref{fig:social_norms_convergence_comparision} (right), all three E2 runs exhibit a rapid rise in convergence by rounds~2--3, then steadily surpass the $8$-point threshold and concentrate in the $9$--$10$ range. 
Across-run variability is noticeably smaller than in E1, where peer-only exchanges leave some runs stalled around medium convergence levels. 
Because $S_t^{\text{conv}}$ directly reflects agreement over proposed norms, this pattern suggests that local peer communication alone does not reliably resolve conflicting norms within our interaction horizon. 
By contrast, the mediation-enabled Summary Agent aggregates and reframes proposals into a shared summary that serves as a common focal point, making it easier for agents to revise their initial positions and reach high, stable convergence.